\documentclass[10pt]{article}

% MoML 2026 short paper (non-archival). Format reproduces the MoML example style
% (log_2022.sty, itself based on Learning on Graphs 2022): US Letter, single column,
% Times 10pt, title rules, a running header, and a first-page venue footer, so the
% build has no external style-file dependency. Build: tectonic paper/moml2026_shortpaper.tex
\usepackage[letterpaper,left=1.5in,right=1.5in,top=1in,bottom=1in]{geometry}
\usepackage{mathptmx}
\usepackage[T1]{fontenc}
\usepackage[utf8]{inputenc}
\usepackage{amsmath,amssymb}
\usepackage{booktabs}
\usepackage{graphicx}
\usepackage{caption}
\usepackage[dvipsnames]{xcolor}
\usepackage{titlesec}
\usepackage{fancyhdr}
\usepackage[numbers,sort&compress]{natbib}
\usepackage[hidelinks]{hyperref}

% MoML: 9pt captions, bold label, sentence case; single-line centered, multi-line flush left.
\captionsetup{font=small,labelfont=bf,justification=raggedright,singlelinecheck=true}
\titlespacing*{\section}{0pt}{7pt}{2pt}
\titlespacing*{\subsection}{0pt}{5pt}{1pt}
\titleformat{\section}{\normalfont\large\bfseries}{\thesection}{0.6em}{}
\titleformat{\subsection}{\normalfont\normalsize\bfseries}{\thesubsection}{0.6em}{}
\setlength{\parindent}{0pt}
\setlength{\parskip}{5.5pt}

% ---- MoML running header (title on pages >= 2) and first-page venue footer ----
\newcommand{\runningtitle}{Know When to Fold: Selective Prediction for Co-Folding}
\fancypagestyle{moml}{%
  \fancyhf{}%
  \renewcommand{\headrulewidth}{0.4pt}%
  \fancyhead[C]{\small\runningtitle}%
  \fancyfoot[C]{\small\thepage}%
}
\fancypagestyle{momlfirst}{%
  \fancyhf{}%
  \renewcommand{\headrulewidth}{0pt}%
  % venue line on top, page number on its own line below, so the wide venue string
  % does not collide with the centered page number.
  \fancyfoot[C]{\footnotesize Massachusetts Institute of Technology -- Molecular Machine Learning Conference 2026.\\[3pt]\normalsize\thepage}%
}
\setlength{\footskip}{30pt}
\pagestyle{moml}

\begin{document}

% ---- MoML title block: rule, title, rule, authors ----
\begin{center}
{\vspace{-1.5em}\rule{\textwidth}{0.8pt}}\\[0.6em]
{\Large\bfseries Know When to Fold: Distribution-Shift-Aware Selective\\ Prediction for Protein--Ligand Co-Folding}\\[0.5em]
{\rule{\textwidth}{0.8pt}}\\[1.1em]
{\bfseries Rikhin Kavuru}\\[0.2em]
{\small Independent Researcher}\\
{\small\texttt{rikhin@virahacks.com}}\\[0.4em]
{\itshape\small MoML 2026 short paper. Non-archival.}
\end{center}
\thispagestyle{momlfirst}

\begin{center}\textbf{Abstract}\end{center}
\vspace{-0.3em}
\begin{quote}
Protein--ligand co-folding models (AlphaFold3, Boltz, Chai) report confidence scores that correlate
with pose accuracy, but a correlation is not a decision rule and it is weakest exactly where drug
discovery operates, on novel pockets and novel chemotypes. We recast co-folding reliability as
selective prediction and build a model-agnostic, training-free layer that turns confidence into a
risk-controlled accept/abstain decision with a finite-sample guarantee. On the released Runs N' Poses
benchmark (five governed co-folding models, 12{,}602 delivered poses) a conformal selective-risk gate
is valid i.i.d.\ yet breaks under training-similarity shift: a marginally compliant gate
under-controls error about twofold on novel strata, and calibrating on familiar ligands then
deploying on novel ones violates the guarantee in every run (realized error 0.55 against a 0.20
target). Group-conditional and weighted conformal keyed on training-set similarity repair per-stratum
coverage, and a combined reliability score cuts the risk-coverage-curve area (AURC) by a quarter to
two-fifths over native confidence, roughly doubling to tripling the predictions retained at a fixed
guarantee. Structural pose-agreement features from the model's own diffusion samples push the AURC
reduction to 38--51\%, and abstention raises a non-circular downstream metric, crystal-contact
recovery (0.90 versus 0.72 on accepted versus rejected poses). The layer ships as a pip-installable
package over frozen model outputs.
\end{quote}

\section{Introduction}
Co-folding models are increasingly used to prioritise molecules and generate starting structures for
downstream design. Recent work establishes that their confidence tracks accuracy and that accuracy
drops on inputs dissimilar to training~\citep{rnp}. Neither gives a practitioner what they need: a
rule for \emph{which predictions to trust}, with a guarantee that survives the distribution shift
they actually face.

We provide that rule. Our contribution is (i) selective prediction and conformal risk control for
co-folding \emph{pose} reliability, a structured 3D output rather than a scalar property. The closest
prior work, CoDrug~\citep{codrug}, conformalizes a scalar molecular property under shift, and we are
not aware of a prior conformal treatment of pose correctness. Further, (ii) a sharp, quantified
demonstration that standard conformal guarantees collapse under novel-pocket and novel-chemotype
shift, using training-set structural and chemical similarity as the shift variable; and (iii) a
shift-robust repair (group-conditional and importance-weighted conformal keyed on that similarity)
plus a combined reliability score, packaged as a training-free layer over frozen AlphaFold3 (AF3),
Boltz, and Chai outputs. We show ordinary probability calibration is not a substitute: it carries no
finite-sample coverage and breaks under the same shift, whereas the conformal variants repair it.

\section{Method}
\textbf{Setup.} The unit is the delivered pose: the top-1 sample per (system, ligand, model) by the
model's own ranking score. The label is $Y=1$ iff the symmetry-corrected ligand root-mean-square
deviation (RMSD) is at most $2$\,\AA. The gate accepts a pose when its confidence $s\ge\tau$ and
abstains otherwise; we report selective risk (error among accepted) against coverage (fraction
accepted).

\textbf{Certified threshold.} We choose $\tau$ by Learn-then-Test with fixed-sequence
testing~\citep{angelopoulos2021ltt} over a pre-specified, data-independent sequence of top-$k$ accept
ranks (smallest first). Each null $H_0(\tau)\!:\,P(\text{error}\mid s\ge\tau)\ge\alpha$ is tested with
the binomial tail p-value $P(\mathrm{Bin}(n_{\text{accept}},\alpha)\le\text{errors})$; the accept set
grows while $H_0$ is rejected at level $\delta$ and stops at the first failure, giving
$P(\text{selective risk}\le\alpha)\ge 1-\delta$, finite-sample and distribution-free. The p-value is
exact for a homogeneous error rate and conservative, hence still valid, for the heterogeneous
Poisson-binomial case by a convex-order argument; conditioning on the accept count is the correct tool
for error-among-accepted.

\textbf{Novelty and shift-robustness.} Training-set similarity (ECFP4 Tanimoto to the nearest training
ligand, pocket similarity, temporal cutoff), shipped pre-computed by Runs N' Poses, is the
covariate-shift variable. It both stratifies (group-conditional / Mondrian conformal~\citep{mondrian}:
a separate $\tau$ per novelty stratum) and weights (weighted conformal~\citep{tibshirani2019}:
reweight familiar-source calibration points by an estimated likelihood ratio to certify on a novel
target without target labels). Ligands with no training analog form their own extreme stratum.

\textbf{Combined score.} Native ranking score is a weak sorter for some models. We fit a
calibration-only combiner (gradient boosting to $P(\text{correct})$) over native confidence, interface
chain-pair ipTM, PoseBusters~\citep{posebusters} physical validity, intra-model ensemble spread across
diffusion samples, ligand difficulty, and cross-model agreement where other models' predictions are
available (a single-model deployment drops this feature, keeping the 22--38\% gain without it).
``Training-free'' here means the co-folding model is never retrained; the combiner is a small model
fit only on the calibration fold, so the layer is free of structure-predictor training, not of all
fitting. Novelty is excluded from the score and enters only through calibration. A 3-way split (fit
combiner / calibrate $\tau$ / test) preserves conformal validity for the learned score: the combiner
never sees the test fold and the threshold is calibrated on out-of-sample combiner scores, so the
learned score carries no combiner-fit leakage (target-clustering across folds is addressed under
Utility).

\section{Results}
Data: released Runs N' Poses predictions consumed as-is (no inference), five governed models
($12{,}602$ delivered poses); a Boltz-2 affinity head brings the raw table to $13{,}535$ and is used
only as an ungoverned comparator. $\alpha=0.20$, $\delta=0.10$ unless noted.

\textbf{Raw novelty gradient.} AF3 delivered-pose correctness falls from $0.88$ on familiar ligands to
$0.44$ on the most novel with-analog stratum. A single global threshold cannot hold a uniform
guarantee across this gradient.

\textbf{Valid i.i.d.} The certifier's finite-sample guarantee, true selective risk $\le\alpha$ with
probability $1-\delta$, is validated on synthetic data where the true risk is known (the gate holds in
$\ge 1-\delta$ of draws). The realized-risk indicator on a finite test fold is a downward-biased proxy
for that event, and we demonstrate the bias rather than assert it: on synthetic data the small-fold
indicator dips below $1-\delta$ while the large-sample (true-risk) indicator meets it. On RNP the
fraction of splits holding, with a Clopper-Pearson interval, reaches the $1-\delta=0.90$ target for
Boltz-1x ($0.90$ $[0.88,0.93]$) and approaches it for Boltz-1 ($0.88$ $[0.85,0.91]$) and AF3 ($0.85$
$[0.81,0.89]$), whose interval sits just below $0.90$; the tight-certifier realized-risk indicator is
the reason it is a downward-biased proxy rather than the guarantee itself, and the certified gate's
mean realized risk is $\le\alpha$ there (AF3 $0.17$ at $28\%$ coverage). Native ranking score certifies
almost nothing for Chai and Protenix ($<5\%$ coverage), a near-vacuous gate that motivates the
combined score below.

\textbf{The break.} AF3's marginal risk ($0.177$) hides severe per-stratum under-control
(Table~\ref{tab:break}). Calibrating on familiar ligands and deploying on novel ones accepts $98\%$ of
poses at realized error $0.55$ with the guarantee holding in $0\%$ of runs. All five governed models
show the same collapse. The break is a property of structural and chemical novelty, not recency: it is
at least as strong along pocket-novelty (AF3 $S_0$ $0.06$ to $S_3$ $0.40$, Chai up to $0.77$) but weak
along a temporal release-date axis, which sharpens what the shift variable must capture.

\begin{table}[t]
\centering
\caption{The exchangeability break: AF3 realized selective risk per novelty stratum ($S_0$ familiar to
$S_4$ no analog) under a single global threshold. The marginal risk of $0.177$ hides the tail.}
\label{tab:break}
\begin{tabular}{lccccc}
\toprule
novelty stratum & $S_0$ & $S_1$ & $S_2$ & $S_3$ & $S_4$ \\
\midrule
realized selective risk & $0.07$ & $0.15$ & $0.16$ & $\mathbf{0.38}$ & $\mathbf{0.43}$ \\
\bottomrule
\end{tabular}
\end{table}

\textbf{The repair.} Group-conditional calibration (a separate $\tau$ per novelty stratum) restores
per-stratum validity: each stratum's accepted error is $\le\alpha$ with probability $1-\delta$
marginally (a simultaneous across-strata statement follows by calibrating each at $\delta/S$ over the
$S$ strata). With native confidence the cost is heavy abstention on the hardest strata (the gate folds
rather than assure falsely). With the combined score, group-conditional calibration instead recovers
usable coverage across the gradient while holding risk $\le\alpha$: AF3 accepts $52\%$ of stratum $S_1$
and $39\%$ of $S_2$ (against $8\%$ and $12\%$ with native confidence) and certifies a small fraction of
the novel $S_3$ that native confidence must abstain on. Only the no-training-analog extreme ($S_4$)
stays uncertifiable, where abstention is the correct answer.

\textbf{Utility.} The combined score dominates native confidence on the risk-coverage curve
(Table~\ref{tab:aurc}). The AURC reduction is significant per model: a paired bootstrap over test poses
gives $90\%$ CIs on $\Delta$(AURC) that exclude zero (AF3 $0.070$ $[0.049,0.092]$, Protenix $0.092$
$[0.068,0.116]$; Boltz-2 omitted for power, $n=933$). Because a target recurs across roughly five
poses, the pooled random split leaks $95\%$ of test targets into calibration; a target-grouped split
and a target-cluster bootstrap give the same conclusion (AF3 $\Delta$(AURC) $[0.056,0.079]$), so the
gain is not an artifact of clustering. At $\alpha=0.2$ the combined gate retains far more predictions at
the same guarantee (AF3 coverage $0.22\to0.71$; Chai $0.01\to0.60$) and unlocks the stringent
$\alpha=0.1$ ($90\%$-correct) operating point that native confidence cannot certify at all for several
models.

\begin{table}[t]
\centering
\caption{Area under the risk-coverage curve (AURC; lower is better), native versus combined score.
The combined score cuts AURC by a quarter to two-fifths, with every per-model $\Delta$(AURC) 90\% CI
excluding zero (five models tested; Boltz-2 omitted for power).}
\label{tab:aurc}
\begin{tabular}{lccc}
\toprule
model & AURC native & AURC combined & $\Delta$ \\
\midrule
AF3      & $0.187$ & $0.125$ & $-33.1\%$ \\
Boltz-1  & $0.234$ & $0.170$ & $-27.5\%$ \\
Boltz-1x & $0.216$ & $0.163$ & $-24.6\%$ \\
Chai-1   & $0.248$ & $0.150$ & $-39.6\%$ \\
Protenix & $0.281$ & $0.175$ & $-37.7\%$ \\
\bottomrule
\end{tabular}
\end{table}

\textbf{Weighted repair, label-free.} Calibrating on familiar ligands and correcting with cross-fitted,
probability-calibrated likelihood-ratio weights over the novelty covariates pulls the realized error
on a moderately-novel target toward $\alpha$ without target labels (AF3 naive $0.27$ to weighted $0.19$
at $0.70$ coverage; same trend for all models, and robust to the weight-model choice). We also
implement an importance-weighted Learn-then-Test gate~\citep{almeida2025} built on a WSR betting
p-value~\citep{waudbysmith2024} that is finite-sample valid conditional on correct weights; on real
co-folding data it abstains, because even the plug-in barely clears $\alpha$, so no certifiable margin
remains. A concept-shift diagnostic explains why: $P(\text{correct}\mid\text{confidence})$ itself moves
between source and target, and the gap grows from the moderate regime ($0.08$--$0.16$) to the extreme
regime ($0.17$--$0.29$), so pure covariate reweighting cannot restore validity on novel pockets.
Weighted conformal is therefore the label-free complement; group-conditional calibration is the
rigorous finite-sample guarantee where stratum labels exist.

\textbf{Baselines: calibration is not conformal.} The combined score has the lowest AURC among the
confidences a practitioner would reach for (AF3 $0.12$ versus native ipTM $0.16$ versus ranking score
$0.19$; same ordering for all models), and PoseBusters validity alone is a poor gate (realized error
$0.26$--$0.34$). The decisive comparison isolates the guarantee from the features: the same combined
score as a Platt/isotonic-calibrated classifier with a fixed threshold versus inside the conformal
layer. Under i.i.d.\ data both control error; under the novelty shift the Platt-calibrated
fixed-threshold gate breaks like naive conformal (realized error $0.22$ above $\alpha$). Isotonic
calibration sits near $\alpha$ ($0.18$) only by abstaining down to $71\%$ coverage, and even then holds
the guarantee in $76\%$ of runs rather than the required $90\%$, since calibration carries no
finite-sample coverage. Only the shift-robust conformal repair (group-conditional, which has no
calibration analogue) restores realized error $\le\alpha$ with a guarantee, by abstaining more on the
novel strata.

\textbf{Downstream payoff.} The gate cleans the delivered pose set a downstream pipeline would carry
forward: base top-1 purity $63$--$70\%$ rises to $82\%$ ($\alpha=0.2$, retaining $86\%$ of correct AF3
poses) or $91$--$94\%$ ($\alpha=0.1$, all models), with interface LDDT-PLI lifting from about $0.72$ to
$0.85$--$0.92$. Purity equals $1$ minus selective risk by construction, so it cannot show the gate buys
anything beyond the guarantee; we add a genuinely different, downstream metric: protein-ligand
interaction-fingerprint recovery (receptor residues within $4.5$\,\AA\ of the ligand versus the crystal
structure), which is only correlated with the $2$\,\AA\ label. Under the gate, accepted poses recover
$0.90$--$0.91$ of the crystal contacts versus $0.71$--$0.76$ for rejected poses, and the
accepted-minus-rejected recall gap 90\% CI excludes zero for all five models (gaps $0.14$--$0.21$).
Abstention routes forward the poses whose interaction patterns a downstream step can trust. A
screening-enrichment harness is implemented; the archival version applies it to a released co-folded
screen~\citep{shen2026vs}, where the governed pose signal enriches actives above docking but the
headline enrichment rides on an ungoverned affinity ranker, so we scope that arm out here.

\textbf{Robustness, task, and label.} The break and the gain do not depend on the $2$\,\AA\ convention:
across thresholds $1.5$--$3.0$\,\AA\ the novel strata always exceed $\alpha$ (AF3 $S_3$ risk
$0.35$--$0.43$ versus $S_0$ $0.05$--$0.09$) and the combined score cuts AURC $31$--$34\%$. On interface
quality (LDDT-PLI $\ge 0.5$) the combined score again dominates native confidence on AURC
($38$--$55\%$ lower). Adding pose-agreement features from the model's own $25$ diffusion samples
(intra-model diversity and cross-model delivered-pose RMSD, symmetry-corrected, no GPU) lowers AURC
further, the single largest drop after ipTM (AF3 $0.123\to0.106$), raising the overall reduction to
$38$--$51\%$ across models. Dropping the threshold and ordering by continuous RMSD, we certify a
continuous gate (mean bounded-RMSD among accepted $\le$ a target with probability $1-\delta$) with a
variance-adaptive WSR betting bound that certifies non-trivial coverage where a worst-case Hoeffding
bound certifies almost none (AF3, $1$\,\AA\ mean-RMSD target: $43\%$ versus $0\%$).

\section{Discussion and limitations}
The guarantee is over the chosen correctness definition and the sampled shift axes, not a universal
notion of trust; the findings are robust across RMSD thresholds $1.5$--$3.0$\,\AA\ and extend to a
certified continuous-RMSD gate. The weighted-conformal guarantee is exact only conditional on correct
importance weights, and under novel-pocket shift there is residual concept shift, so we keep
group-conditional as the rigorous finite-sample guarantee and scope weighted conformal as the
label-free complement. A cross-dataset test on FoldBench recovers the withheld interface-ipTM by
regenerating the Protenix predictions and self-scoring ligand-RMSD against the deposited assemblies
($436$ targets, $384$ train-similar, $52$ low-homology; feature and label from one run). This compares
gates, not benchmarks: our self-scored Protenix top-1 success is $0.40$ against FoldBench's released
$0.57$, a gap from the scoring convention, model version, and MSA source that does not affect the
transfer because feature and label are self-consistent. The frozen Runs N' Poses interface-ipTM gate
transfers with a risk-coverage AURC of $0.380$, ahead of the matched ranking-score control at $0.454$,
and the same threshold accepts only $5\%$ of FoldBench against $26\%$ at home, so the coverage the gate
can certify collapses on the low-homology second dataset. On the extreme novel-chemotype regime (base
correctness $<55\%$) no gate can certify a high-coverage low-error set; the layer's value there is
principled abstention over a naive gate's false confidence. Code, calibration tables, and a
one-command reproduction over the released predictions are provided in the accompanying repository
(public release on acceptance).

\small
\bibliographystyle{plainnat}
\begin{thebibliography}{9}
\bibitem{angelopoulos2021ltt} A. Angelopoulos, S. Bates, E. Cand\`es, M. Jordan, L. Lei. Learn then Test: Calibrating Predictive Algorithms to Achieve Risk Control. 2021.
\bibitem{tibshirani2019} R. Tibshirani, R. Foygel Barber, E. Cand\`es, A. Ramdas. Conformal Prediction Under Covariate Shift. NeurIPS 2019.
\bibitem{codrug} K. Nguyen et al. CoDrug: Conformal Drug Property Prediction with Density Estimation under Covariate Shift. NeurIPS 2023.
\bibitem{mondrian} V. Vovk. Conditional Validity of Inductive Conformal Predictors (Mondrian). Machine Learning 2013.
\bibitem{waudbysmith2024} I. Waudby-Smith, A. Ramdas. Estimating Means of Bounded Random Variables by Betting. JRSS-B 2024.
\bibitem{almeida2025} Almeida et al. High-probability, importance-weighted risk control. 2025.
\bibitem{rnp} A. \v{S}krinjar et al. Have protein--ligand co-folding methods moved beyond memorisation? (Runs N' Poses) 2025.
\bibitem{shen2026vs} C. Shen et al. Unlocking the application potential of AlphaFold3-like co-folding for virtual screening. Chem. Sci. 2025.
\bibitem{posebusters} M. Buttenschoen, G. Morris, C. Deane. PoseBusters. Chem. Sci. 2024.
\end{thebibliography}

\end{document}
